% =============================================================
%  CS102-F26 Digital Design and Computer Architecture I
%  Lecture 21 -- The Stepper, the Data Path and the Controller (Unit 21)
%  Follows Sections 6.6.2 and 6.6.3 of
%  A Student's Computer Architecture Book I.
%
%  Build the slides and the print handout:  npm run materials -- lecture-21
% =============================================================
\makeatletter
\def\input@path{{templates/latex/}{../}}
\makeatother

\documentclass[10pt]{wosm-slides}
\usepackage{cs102-lectures}[2026/10/20]

\title[Lecture 21 \textperiodcentered\ Stepper, Data Path, Controller]{The Stepper,
  the Data Path and the Controller}
\subtitle{Lecture 21: how the DX8 runs an instruction, step by step}

% ---------- Helpers for this lecture --------------------------
\tikzset{
  part/.style={draw=wosm-text, line width=0.6pt, fill=wosm-bg, align=center,
    font=\scriptsize, minimum width=1.5cm, minimum height=0.75cm, inner sep=2pt},
  group/.style={draw=wosm-muted, line width=0.5pt, dashed},
  grouplabel/.style={font=\scriptsize, text=wosm-muted},
  hot/.style={draw=wosm-brick, line width=1.5pt},
}
\tikzset{gates/.append style={
  and gate/.append style={logic gate input sep=0.3cm}}}

% The DX8 of the book's Figures 6.8 to 6.11. The argument makes some paths hot,
% e.g. \datapath{ip-ram/.style=hot, st-ctl/.style=hot}.
\newcommand*\datapath[1]{%
  \begin{tikzpicture}[x=0.92cm, y=0.92cm,
      ip-ram/.style={}, ram-rf/.style={}, rf-ir/.style={}, rf-ab/.style={},
      ra-alu/.style={}, rb-alu/.style={}, ir-alu/.style={}, st-ctl/.style={},
      ctl-alu/.style={}, alu-out/.style={}, alu-flag/.style={}, out-rf/.style={},
      flag-rf/.style={}, ir-ctl/.style={}, ctl-ip/.style={}, #1]
    % groups
    \draw[group] (1.65,-1.85) rectangle (3.55,1.75);
    \node[grouplabel, anchor=north] at (2.6,-1.85) {memory unit};
    \draw[group] (4.0,-2.85) rectangle (11.45,2.4);
    \node[grouplabel, anchor=south east] at (11.45,-2.85) {CPU};
    % parts
    \node[part, minimum width=1.4cm] (in) at (0,0) {input unit\\ + buttons};
    \node[part, minimum width=1.4cm, minimum height=0.8cm] (ram) at (2.6,1.0) {RAM};
    \node[part, minimum width=1.4cm] (ip) at (2.6,-1.05) {instruction\\ pointer (IP)};
    \node[part, minimum height=0.8cm] (rf) at (5.3,1.0) {register file\\ R0--R7};
    \node[part] (ir) at (5.3,-0.3) {instruction\\ register (IR)};
    \node[part, minimum width=0.9cm, minimum height=0.5cm] (ra) at (7.2,1.55) {regA};
    \node[part, minimum width=0.9cm, minimum height=0.5cm] (rb) at (8.4,1.55) {regB};
    \node[part, minimum width=2.2cm, minimum height=0.9cm] (alu) at (7.8,0.1) {ALU};
    \node[part, minimum width=1.2cm, minimum height=0.45cm] (out) at (7.6,-1.0) {regOUT};
    \node[part, minimum width=1.2cm, minimum height=0.45cm] (flag) at (7.6,-1.7) {regFLAG};
    \node[part, minimum width=1.4cm] (st) at (10.4,1.55) {stepper};
    \node[part, minimum width=1.4cm] (ctl) at (10.4,0.1) {controller};
    % input
    \draw[wire] (in.east) -- (1.1,0);
    \draw[arrow] (1.1,0) |- (ram.west);
    \draw[arrow] (1.1,0) |- (ip.west);
    \fill[wosm-text] (1.1,0) circle (1.5pt);
    % paths
    \draw[arrow, ip-ram] (ip) -- (ram);
    \draw[arrow, {Stealth[length=5pt]}-{Stealth[length=5pt]}, ram-rf]
      (ram.east |- 0,1.15) -- (rf.west |- 0,1.15);
    \draw[arrow, rf-ir] (rf) -- (ir);
    \draw[wire, rf-ab] (rf.east) -- ++(0.25,0) |- (8.4,2.1);
    \draw[arrow, rf-ab] (7.2,2.1) -- (ra);
    \draw[arrow, rf-ab] (8.4,2.1) -- (rb);
    \fill[wosm-text] (7.2,2.1) circle (1.5pt);
    \draw[arrow, ra-alu] (ra) -- (ra |- alu.north);
    \draw[arrow, rb-alu] (rb) -- (rb |- alu.north);
    \draw[arrow, ir-alu] (ir.east) -- (ir.east -| alu.west);
    \draw[arrow, st-ctl] (st) -- (ctl);
    \draw[arrow, ctl-alu] (ctl) -- (alu);
    \draw[arrow, alu-out] (alu.south -| out) -- (out);
    \draw[arrow, alu-flag] (8.7,-0.39) |- (flag.east);
    \draw[wire, flag-rf] (flag.west) -| (4.25,-1.0);
    \draw[arrow, out-rf] (out.west) -| (4.25,0.85) -- (rf.west |- 0,0.85);
    \fill[wosm-text] (4.25,-1.0) circle (1.5pt);
    \draw[arrow, ir-ctl] (ir.south) |- (10.1,-2.35) -- (10.1,0 |- ctl.south);
    \draw[arrow, ctl-ip] (10.75,0 |- ctl.south) |- (2.6,-2.6) -- (ip.south);
  \end{tikzpicture}}

% =============================================================
\begin{document}

\maketitle

\begin{frame}{Today}
  \tableofcontents
\end{frame}

% =============================================================
\section{The stepper}

\begin{frame}{The stepper}
  \begin{columns}[c, totalwidth=\textwidth]
    \column{0.5\textwidth}
      \centering
      \begin{tikzpicture}
        \node[part, minimum width=1.8cm, minimum height=1.5cm] (c) at (0,0)
          {3-bit\\ synchronous\\ up counter};
        \node[part, minimum width=1.1cm, minimum height=3.0cm] (d) at (2.4,0)
          {3-to-8\\ decoder};
        \foreach \i/\y in {2/0.4, 1/0, 0/-0.4} {
          \draw[arrow] (c.east |- 0,\y) -- node[above, font=\tiny] {$y_\i$}
            (d.west |- 0,\y);
        }
        \foreach \k in {1,...,8} {
          \pgfmathsetmacro\y{1.225-0.35*(\k-1)}
          \draw[arrow] (d.east |- 0,\y) -- ++(0.5,0) node[right, font=\tiny] {$s_\k$};
        }
        \draw[arrow] ($(c.west)+(-0.6,0.35)$) node[left, font=\scriptsize] {clk} --
          ($(c.west)+(0,0.35)$);
        \draw[arrow] ($(c.west)+(-0.6,-0.35)$) node[left, font=\scriptsize] {reset} --
          ($(c.west)+(0,-0.35)$);
      \end{tikzpicture}\par\smallskip
      \wosmlabel{After the book's Figure 6.5}
    \column{0.46\textwidth}
      \begin{itemize}
        \item The stepper turns on one output at a time, in a fixed order, at each
          clock edge: its outputs are \emph{one-hot} \unitref{8}.
        \item Built from parts we know: a 3-bit synchronous up counter
          \unitref{13} and a 3-to-8 decoder \unitref{8}.
        \item Its one input besides the clock is \emph{reset}: the counter goes to
          000, so the stepper goes to step 1.
      \end{itemize}
  \end{columns}
\end{frame}

\begin{frame}{Steps and stages}
  \begin{center}
    \begin{tikzpicture}[x=0.5cm]
      \foreach \t in {0,2,...,16} {\timingedge{\t}{8}}
      \timingwave{0}{clk}{1,0,1,0,1,0,1,0,1,0,1,0,1,0,1,0,1,0}
      \timingwave{1}{$s_1$}{1,1,0,0,0,0,0,0,0,0,0,0,0,0,0,0,1,1}
      \timingwave{2}{$s_2$}{0,0,1,1,0,0,0,0,0,0,0,0,0,0,0,0,0,0}
      \timingwave{3}{$s_3$}{0,0,0,0,1,1,0,0,0,0,0,0,0,0,0,0,0,0}
      \timingwave{4}{$s_4$}{0,0,0,0,0,0,1,1,0,0,0,0,0,0,0,0,0,0}
      \timingwave{5}{$s_5$}{0,0,0,0,0,0,0,0,1,1,0,0,0,0,0,0,0,0}
      \timingwave{6}{$s_6$}{0,0,0,0,0,0,0,0,0,0,1,1,0,0,0,0,0,0}
      \timingwave{7}{$s_7$}{0,0,0,0,0,0,0,0,0,0,0,0,1,1,0,0,0,0}
      \timingwave{8}{$s_8$}{0,0,0,0,0,0,0,0,0,0,0,0,0,0,1,1,0,0}
      \foreach \x/\t in {2/fetch, 6/decode, 10/execute, 14/load} {
        \node[font=\scriptsize, text=wosm-brick] at (\x,0.75) {\t};
      }
      \foreach \x in {4,8,12} {\draw[wosm-brick, line width=0.4pt] (\x,0.95) --
        (\x,0.55);}
    \end{tikzpicture}
  \end{center}

  \begin{itemize}
    \item Two steps per stage. After $s_8$ comes $s_1$ again, and the next
      instruction.
  \end{itemize}
\end{frame}

% =============================================================
\section{The data path}

\begin{frame}{Which wires are hot?}
  \begin{itemize}
    \item The \emph{data path} is the set of wires along which data moves between
      the parts of the CPU and the memory.
    \item In each stage, only some of them carry data that matters. The book draws
      them in orange for each stage (Figures 6.8 to 6.11); here they are in red.
    \item To make data go from part A to part B, the controller \alert{read-enables}
      A, \alert{write-enables} B, and sets the select bits of the multiplexers on the
      way.
  \end{itemize}

  \begin{remark}
    Every stage below also uses the stepper's output, so that the controller knows
    which step it is in.
  \end{remark}
\end{frame}

\begin{frame}{Fetch}
  \begin{center}
    \datapath{ip-ram/.style=hot, ram-rf/.style=hot, st-ctl/.style=hot,
      ctl-ip/.style=hot}
  \end{center}
  \begin{itemize}
    \item The IP gives the RAM its address; the instruction comes from the RAM into
      the CPU. The controller also drives the IP.
  \end{itemize}
\end{frame}

\begin{frame}{Decode}
  \begin{center}
    \datapath{rf-ab/.style=hot, rf-ir/.style=hot, ir-alu/.style=hot,
      ctl-alu/.style=hot, ir-ctl/.style=hot, st-ctl/.style=hot}
  \end{center}
  \begin{itemize}
    \item The instruction is in the IR. Its arguments go to regA and regB, its
      operation, 4 bits, to the ALU; the controller reads it too.
  \end{itemize}
\end{frame}

\begin{frame}{Execute}
  \begin{center}
    \datapath{ra-alu/.style=hot, rb-alu/.style=hot, ir-alu/.style=hot,
      ctl-alu/.style=hot, alu-out/.style=hot, alu-flag/.style=hot,
      ir-ctl/.style=hot, st-ctl/.style=hot}
  \end{center}
  \begin{itemize}
    \item For an ALU instruction: regA and regB into the ALU, which computes what the
      IR asks for. The result goes to regOUT, the flags to regFLAG.
  \end{itemize}
\end{frame}

\begin{frame}{Load}
  \begin{center}
    \datapath{out-rf/.style=hot, flag-rf/.style=hot, ram-rf/.style=hot,
      ip-ram/.style=hot, ctl-ip/.style=hot, st-ctl/.style=hot}
  \end{center}
  \begin{itemize}
    \item regOUT goes to the register file; or a register goes to the RAM, at the
      address in the instruction. The controller moves the IP on to the next
      instruction.
  \end{itemize}
\end{frame}

% =============================================================
\section{The controller}

\begin{frame}{What the controller does}
  \begin{columns}[c, totalwidth=\textwidth]
    \column{0.49\textwidth}
      \centering
      \begin{tikzpicture}
        \node[part, minimum width=1.8cm, minimum height=2.4cm] (c) {controller};
        \foreach \t/\y in {{steps $s_1 \ldots s_8$}/0.8, {$b_{15}$ (ALU or not)}/0.27,
            {opcode $b_{14:11}$}/-0.27, {flags}/-0.8} {
          \draw[arrow] ($(c.west)+(-0.5,\y)$) node[left, font=\scriptsize] {\t} --
            ($(c.west)+(0,\y)$);
        }
        \foreach \t/\y in {{read enables}/0.8, {write enables}/0.27,
            {select bits}/-0.27, {setIP, \ldots}/-0.8} {
          \draw[arrow] ($(c.east)+(0,\y)$) -- ++(0.5,0) node[right, font=\scriptsize]
            {\t};
        }
      \end{tikzpicture}
    \column{0.47\textwidth}
      \begin{itemize}
        \item Inputs: which step it is (from the stepper), the instruction (from the
          IR), and the ALU's flags.
        \item Outputs: the enable and select bits of every other part, for this step
          of this instruction.
        \item It is \alert{combinational}: AND and OR gates, as in Lectures 6 to 8.
          The sequence in time comes from the stepper.
      \end{itemize}
  \end{columns}
\end{frame}

\begin{frame}{Recognising the instruction}
  \begin{columns}[c, totalwidth=\textwidth]
    \column{0.42\textwidth}
      \centering
      \begin{tikzpicture}[gates]
        \node[and gate, inputs={iiiin}] (g) at (0,0) {};
        \foreach \i/\b in {1/15, 2/14, 3/13, 4/12, 5/11} {
          \draw[wire] (g.input \i) -- ++(-0.5,0) node[left, font=\scriptsize]
            {$b_{\b}$};
        }
        \draw[wire] (g.output) -- ++(0.4,0) node[right, font=\scriptsize]
          {\texttt{store}};
      \end{tikzpicture}\par\smallskip
      \wosmlabel{1 exactly for 0\,0001}
    \column{0.54\textwidth}
      \begin{itemize}
        \item One AND gate per instruction: its inputs are $b_{15}$ and the opcode,
          with a bubble on every bit that must be 0. It is a decoder \unitref{8},
          one output at a time.
        \item The book's controller recognises \texttt{move}, \texttt{arg1},
          \texttt{arg2}, \texttt{store} and the four jumps this way.
        \item The ALU instructions need no such gate: their opcode goes straight to
          the ALU's multiplexer \unitref{18}.
      \end{itemize}
  \end{columns}
\end{frame}

\begin{frame}{Some control signals}
  \begin{center}
    \begin{tabular}{@{}ll@{}}
      \toprule
      Signal & In the book's controller (Figure 6.12) \\
      \midrule
      exec & $s_5 + s_6$: the execute stage \\
      load & $s_7 + s_8$: the load stage \\
      RAM\_WE & \texttt{store} $\cdot$ load: the RAM is written only in a
        \texttt{store}'s load stage \\
      load\_aluOUT & $b_{15} \cdot$ load: an ALU result goes back in the load stage \\
      setIP & \texttt{jmp} $+$ \texttt{jmpgt}$\cdot$gt $+$ \texttt{jmplt}$\cdot$lt
        $+$ \texttt{jmpeq}$\cdot$eq \\
      regFileWE & 1 in step $s_1$ (and others): the instruction comes from the RAM \\
      \bottomrule
    \end{tabular}
  \end{center}

  \begin{itemize}
    \item setIP is 1 when the IP must take the jump address \texttt{val}: always for
      \texttt{jmp}, and for a conditional jump only if its flag is set \unitref{15}.
  \end{itemize}
\end{frame}

\begin{frame}{Try it}
  \begin{enumerate}
    \item In step $s_6$, which stage is the DX8 in? Which step follows $s_8$?\par
      \uncover<2->{\small Execute. Then $s_1$: the counter wraps around
        \unitref{13}.}
    \item Draw the AND gate that recognises \texttt{arg1} (0\,0011).\par
      \uncover<3->{\small Inputs $\overline{b_{15}}$, $\overline{b_{14}}$,
        $\overline{b_{13}}$, $b_{12}$, $b_{11}$.}
    \item When is RAM\_WE $= 1$?\par
      \uncover<4->{\small In steps $s_7$ and $s_8$ of a \texttt{store}
        instruction, and never otherwise.}
    \item The IR holds a \texttt{jmpeq} and eq $= 0$. What is setIP?\par
      \uncover<5->{\small 0: no jump, the program goes on to the next instruction.}
    \item Which paths are hot while an \texttt{add} is executed?\par
      \uncover<6->{\small regA, regB, the IR and the controller into the ALU; the
        ALU into regOUT and regFLAG; the IR into the controller.}
  \end{enumerate}
\end{frame}

% =============================================================
\section{Summary}

\begin{frame}{Summary}
  \begin{itemize}
    \item The stepper, a 3-bit counter and a 3-to-8 decoder, gives one-hot steps
      $s_1$ to $s_8$: two for each of fetch, decode, execute and load.
    \item The data path: the wires along which data moves. Each stage uses a
      different set of them.
    \item The controller is combinational. From the step, the instruction and the
      flags, it sets the read enables, write enables and select bits.
    \item Each instruction is recognised by one AND gate on $b_{15}$ and the opcode.
  \end{itemize}
\end{frame}

\begin{frame}{Before the next lecture}
  \begin{itemize}
    \item Read Sections~6.6.2 and 6.6.3 of the book.
    \item In Logisim-evolution, build the stepper from your counter and a decoder, and
      watch $s_1$ to $s_8$ in a timing diagram.
    \item Build the instruction-recognising AND gates and the setIP logic of the
      controller.
  \end{itemize}

  \begin{block}{Next lecture \unitref{22}}
    Putting it all together: building and programming the DX8 computer.
  \end{block}
\end{frame}

\end{document}
